% A self-contained article for one pull request.
\documentclass[11pt]{article}
\usepackage[T1]{fontenc}
\usepackage{lmodern}
\usepackage{microtype}
\usepackage[margin=1in]{geometry}
\usepackage{amsmath,amssymb,amsthm,mathtools}
\usepackage{algorithm}
\usepackage{algpseudocode}
\usepackage{booktabs}
\usepackage{tabularx}
\usepackage[numbers,sort&compress]{natbib}
\usepackage[hidelinks]{hyperref}
\usepackage[capitalise,noabbrev]{cleveref}

\theoremstyle{plain}
\newtheorem{proposition}{Proposition}
\theoremstyle{definition}
\newtheorem{definition}{Definition}
\theoremstyle{remark}
\newtheorem{remark}{Remark}

\newcommand{\code}[1]{\texttt{#1}}

% The read-only check of the corpus instance (2026-11-13, on the tree with the landing log).
\newcommand{\corpusitems}{400{,}274}
\newcommand{\corpusunits}{1{,}369{,}944}
\newcommand{\corpusrows}{1{,}338{,}585}
\newcommand{\corpusentries}{389{,}745}
\newcommand{\lexsec}{17.4}
\newcommand{\lexload}{2.84}
\newcommand{\ntests}{62}

\title{Reset, Destroy and a Shared Cache: the Lifecycle of a Knowledge-Store Instance}
\author{Alberto Serrano-Calva}
\date{\today}

\begin{document}
\maketitle

\begin{abstract}
A knowledge store that keeps the verbatim bytes of everything it ingests can treat every index over
those bytes as a materialized view. We use that to give an instance a lifecycle with two operations
whose meaning is fixed in code. A reset classifies each of the instance's stores as a ledger, an
owner's declaration, a derived store, model output, or operational state, clears the derived stores
and rebuilds them in a fixed order from the raw store and the ledgers, and keeps everything else row
for row. A destroy removes a named instance and leaves its open content in a cache shared by every
instance of the machine, addressed by content, so that another instance can take a raw object by its
digest, a vector by its embedder and text hash, and a downloaded source stream by its manifest
without a trust decision beyond the address. A purge reaches the cache, so the cache is no way
around erasure. The catalog's file also holds a hash-chained log of the catalog's past: a reset
keeps it whole and appends nothing to it, and refuses before writing anything when the catalog's
first open would write the log's genesis without the instance's secret; a destroy removes the log
with the catalog and lets none of it into the cache. Two gaps in the store had to close first: the
chat log could not be rebuilt because nothing recorded which raw objects were transcripts, and a
lexical index made by an older tokenizer refused every landing while the job that would repair it
ran only after a landing. We add a retention ledger for the first and a repair inside the one
writer for the second. \ntests{} tests on
fixture palaces check that a reset keeps every ledger, the log's rows and chain heads among them,
converges to the state before for each derived store, is idempotent, calls no embedder while the
atoms are held, and finishes after an interruption at any step, and that a destroy removes only
the instance's own state and leaves exactly the cache's classes.
\end{abstract}

\section{Introduction}

The store this article concerns keeps every ingested input verbatim in a content-addressed raw
store and lands it through one writer that records the item in a catalog, a version chain and an
attestation ledger, then derives its lexical units and its vectors. Several instances of the store
run on one machine: the owner's default instance, a code instance that indexes repositories, and a
corpus instance holding \corpusitems{} open-text items. Two operations on an instance had no single
meaning. The reset at the starting commit deleted the raw store, the catalog, the attestations, the
verdicts and most derived stores, so it lost the only copy of transcripts that their source
application had already pruned, and it left the lexical index and the relation edges behind,
pointing at items that no longer existed. There was no destroy.

The owner's decisions fix the meanings: a reset clears what is derived and rebuilds it from raw, so
it loses nothing; a destroy removes an instance and keeps only what another instance can reuse as a
cache; the corpus instance lives at its own root outside every checkout. Two standing rules bind the
build: a change replaces what it changes, with no switch back to the old behaviour, and what an
operation computes is defined once and is the same in every instance.

The contributions are: a class table of every store an instance writes (\cref{sec:inventory}); a
reset that rebuilds the derived stores in an order where each step reads only ledgers and earlier
steps, with a marker that refuses every writer that starts before it finishes (\cref{sec:reset}); a retention
ledger that makes the chat log a fold of retained transcripts and a lexical repair inside the
writer (\cref{sec:gaps}); a shared, content-addressed cache with an allow-list of open content and
erasure that reaches it (\cref{sec:destroy}); and tests of each property (\cref{sec:eval}).

\section{Background}

Helland argues that once storage is cheap the log of what happened is the source of truth and every
database is a cache of part of it \citep{helland2015immutability}. A derived store is then a
materialized view, and its maintenance a view-maintenance problem \citep{gupta1995views}; a full
recomputation is the simplest correct maintenance. Build systems face the same question of what to
rebuild and in which order, and separate the scheduler from the rebuilding strategy; a rebuild that
keeps a memo of a pure function's results, indexed by a hash of its input, and fetches a result
from it instead of recomputing is their constructive trace, and a memo shared between instances is
their cloud build \citep{mokhov2018build}. Content addressing makes a memo shareable: Venti
names every block by its hash so that identical data is stored once and any holder of the address
can verify what it receives \citep{quinlan2002venti}, and Nix builds a shared store and binary
caches on the same principle, with garbage collection of what nothing references
\citep{dolstra2004nix}. Erasure across derived state is harder: a deletion in the source must reach
every artifact computed from it, which Cao and Yang address for learned models
\citep{cao2015unlearning}; here the derived artifacts are indexes and a cache, and the problem is to
name every one.

\section{The inventory}\label{sec:inventory}

\begin{definition}
A store is a \emph{ledger} when it records something no other store records (the raw objects, the
catalog, the version chains, the attestations, the commit ledger, the reference edges, the run
ledgers). It is a \emph{declaration} when the owner wrote it and no code can compute it (verdicts,
authored supersessions). It is \emph{derived} when a function without a model computes it from
ledgers. It is \emph{kept} when it is model output or otherwise cannot be computed again exactly. It
is \emph{operational} when it is state of the running system and not content.
\end{definition}

The classes are one table in code (\code{ops/lifecycle/inventory.py}), each entry with the path it
lies at under a configuration and, for a derived store, its rebuild step. The reset's two older
lists, its targets and its guard, are replaced by it. A test exercises every lane on a fixture and
fails if any entry of the data directory is unclassified.

Three classifications are choices with consequences. The catalog is a ledger, and its file holds
a second one: the landing log, an append-only, hash-chained record of which digest each path held
when, written in the transaction of each catalog row, so that the catalog's rows are the log's
fold. A transcript or a dialogue export the catalog was landed from may be gone, so the log can be
computed from nothing else; the reset keeps the file whole. The reference edges are a ledger: their
rows carry no extractor version, so a re-projection would not give them back. The code lane's
membership store is kept: it holds the code atoms' embedder ledger beside their occupancy, and
clearing it would forget which code atoms are held under which embedder.

\section{Reset}\label{sec:reset}

\begin{algorithm}[t]
\caption{Reset of an instance $I$ with derived stores $D_1,\dots,D_k$ in rebuild order.}
\label{alg:reset}
\begin{algorithmic}[1]
\Require no live supervisor of $I$; the supervisor lock of $I$ acquired; embedder geometry equals the plane's; every ledgered commit present; the landing log's secret present if the catalog holds rows and no log
\State $F \gets \{D_j : D_j \text{ present}\} \cup \text{stores named by a pending marker}$
\State $H \gets$ the plane's current (path, digest) pairs, or those a pending marker holds
\State write marker $\{F, H\}$ atomically
\State seed the retention ledger from the chat log if it predates the ledger
\ForAll{$D \in F$} \State clear $D$ by its own method (index reset, plane relations, authored rows, or file)
\EndFor
\ForAll{$D \in F$ in step order} \State rebuild $D$ from ledgers and $D_1..D_{j-1}$, writing no ledger
\EndFor
\State remove the marker
\end{algorithmic}
\end{algorithm}

\Cref{alg:reset} gives the reset. The steps are: the lexical index (reset, then the lexical
backfill from the catalog and raw); the plane's occupancies and supersession log, whose atoms are
kept, re-occupied for exactly the pairs $H$, each from raw at its own digest even where the
catalog's head has moved on, so that items a lazy instance has not yet embedded stay with its
embed job; the code and agent observations, re-projected with the attestor left out; the authored
relation edges, parsed from each note's raw bytes by its catalog digest instead of its vault file;
the chat log, replayed from the retention ledger; the chat events; the causal edges; the verdict
dispositions, by applying every verdict in sequence order.

\begin{proposition}[No embedder call]\label{prop:noembed}
Let the plane's atom table map $\mathrm{id}(t) = (\ell, \mathrm{sha256}(t))$, where $\ell$ is the
atom's layer (prose, for a note's chunk), to the embedder's vector of text $t$ under the recorded
geometry $g$. If a reset keeps the table and the configured
geometry equals $g$, the rebuild calls the embedder only for texts whose atom the table lacks.
\end{proposition}
\begin{proof}[Sketch]
The drain embeds exactly the atoms \code{known\_atoms} does not report, and the table reports every
atom it holds. Keeping it loses no derived state, since each row is a value of one function at one
argument. The reset refuses a changed geometry, so no row of another function is consulted.
\end{proof}

\begin{proposition}[Convergence]\label{prop:conv}
For each derived store except the causal edges, the head after a reset equals the head before it,
wall-clock bookmarks excluded, provided every derivation is deterministic in its inputs and the
inputs are kept.
\end{proposition}
The head is what a reader reads: current units and postings, current occupancies with their atoms'
vectors, every row of the chat log, events, dispositions, observations and authored edges. The
version-keyed stores are rebuilt for current versions only; the landing log names every version
landed since its chain began, so a replay of it could rebuild the rest, which is left as future
work. The causal edges resolve commits against the
commit ledger as it stands at each pass, so they equal a fresh integration over the present inputs;
a second reset gives the same edges.

\paragraph{The landing log.} Two steps open the catalog (the lexical index and the authored
edges), and neither writes it: each reads its active rows. The plane's step reads no row of it,
since it rebuilds from raw what the plane held, and does not open it. The one write a reset can
cause there is the catalog's own. The first open of a catalog that holds rows and no log writes the
log's genesis, keyed by the instance's secret, as any first open does; without the secret that
open refuses, but only after creating the log's two empty tables. So the reset reads the catalog
read only before it writes its marker and refuses, naming the secret's item, rather than stopping
after the clearing. With the log in the file the reset needs no secret, and the log's every row
and every chain's head hash are what they were.

The marker makes the reset recoverable: every clear and every step is idempotent, the marker names
every store found, and a second confirmed reset clears and rebuilds them all, including a store
whose file the first had already removed. While the marker exists the daemon's start and every
command that lands or writes a derived store refuse. Each reads the marker when it starts, so it
does not exclude a writer already running; the lock does, for the writers that take it, and the
on-demand writers that take none are not excluded, which is left open. A rerun finishes only if no
step fails the same way every time, so whatever would stop a step that way is found before the
marker is written or tolerated by the step: a derived store outside the data directory, and a
commit the commit ledger holds that the repository has pruned, refuse the reset with nothing
written; a retained transcript whose raw object a purge erased is skipped by the replay and
counted, and so is a version the plane held that a purge erased while the reset was pending, so
the erasure stands.

\section{Two gaps closed}\label{sec:gaps}

\paragraph{The chat log is a fold.} The chat sensor stored each transcript in the raw store and
extracted it only when the raw store reported the bytes new, and it recorded nothing else. With
raw kept, a rebuild would see every transcript as unchanged and extract nothing, and a pruned
transcript would never be read. The order matters too: the first write of a turn stands and a
session that later shows a candidate secret is held at its state before it, so the chat log is the
left fold of retained transcripts in retention order. We add an append-only retention ledger
\code{(seq, digest, file, at)}, written at retention before extraction; its row is now the sensor's
change signal. That closes the gap between the raw write and the row; an extraction cut short after
the row is skipped as unchanged by the next sync, and a reset's replay of the ledger recovers it.
For an instance
whose chat log predates the ledger, the ledger is seeded from the chat log once: one row per
distinct transcript digest, ordered within a session by the lowest turn it wrote. A digest that
wrote no row changes nothing in a replay, so the seed's replay gives back the chat log.

\paragraph{A stale index blocks every landing.} The lexical index records its tokenizer and
refuses to land or score under another. The job that re-derives it was enqueued only after a landing
succeeded, so a tokenizer change left the code instance's repository sync failing on every run
until the index was re-derived by hand (2,021 items in 11~s). The writer now checks the index before
it prices a plan and, when stale, resets it and lands every active item again before the plan. A
test reproduces the stuck state and shows the sync landing.

\section{Destroy and the shared cache}\label{sec:destroy}

A destroy applies to a named instance only and is a dry run unless confirmed; the dry run takes no
lock and writes no file, reading each SQLite file at rest (no \code{-wal}, or an empty one) with
\code{immutable=1}, since a read-only connection to a WAL database creates its \code{-wal} and
\code{-shm} files, and asking the run ledger whether a run is live without opening it, so it
creates no root that does not exist. A destroy reads the catalog read only and removes it whole,
its landing log with it, so it needs no secret. It refuses a live run, a held lock, an installed
agent, a root inside a git working tree, a root that is the home root, holds the cache, or
contains a symbolic link out of itself, a source file the cache would take that is a link, and a
cache volume without room. Since it removes the whole root, it also refuses a root
that holds anything but the instance's own state: a top-level entry no namespace of the instance
names, an unclassified entry of its data directory, a git repository below it, or another declared
instance's root. It fills the cache, then removes the root. A fill refuses unless the configured
embedder is the geometry the plane recorded, so a vector is filed only under the key of the
embedder that made it; its plan counts atoms, and the fill copies vectors in batches.

The cache lives at one place under the home root, with namespaces for raw objects by digest, vectors
by key and atom, and source files with a manifest of size and SHA-256. A vector's key hashes the
embedder's model, document template, dimension and serving backend. An entry carries no path, title
or class. A lookup records the day the entry was taken, for retention, and nothing finer.

\begin{definition}
An item is \emph{open} when its catalog path lies in an open-corpus source's space
(\code{corpus/<name>/}) or a repository source's (\code{code/<name>/}), or its class is curated. A
raw object enters the cache when a catalog row at an open path names its digest; a vector enters
when an open item's occupancy holds its atom; a source file enters when it is a named part, manifest
or metadata file of an open-corpus stream and a plain file of its source: not a symbolic link, not
a file with another hard link, not one that resolves out of the source. Nothing else enters: no
digest a purge recorded, no row of the catalog or of its landing log, and no part of the log's
secret.
\end{definition}

\begin{remark}
Sharing an atom that both an open item and a private note hold reveals nothing beyond the open
source: the entry says only that some open item held this text, which its source publishes. An
entry whose existence said more than the open sources would break the argument; the allow-list
takes open items alone.
\end{remark}

Three readers take entries back: the vector drain looks up an atom under its own key before calling
the embedder, discarding a vector of another width; the raw store, on a miss, takes an object whose
bytes hash to the digest asked for and deletes an entry that does not; and the corpus runner copies
a cached stream in, file by file against the manifest, before its edge reader resumes. A purge,
before deleting a raw object, derives its chunks as the landing did, removes the object and every
atom's vector under every key from the cache, and records the digest, so no fill restores it. The
atoms are deleted rather than recorded, since an atom's text may be held by an open item; the cost
of deleting a shared one is one embedding. Entries not taken for 180 days are removed by a prune.

\section{Evaluation}\label{sec:eval}

\paragraph{Tests.} \ntests{} tests run on fixture palaces that exercise every lane: three notes with
an authored and an interpreted relation edge and a superseded version, three open-corpus items, a
repository landed lexically and read by the code sensor, four transcripts (one grown, one held by
the secret guard, one with a commit that the integrator resolves), signed verdicts, and an
attestation trail. With the transcript directory emptied, a confirmed reset leaves the raw files
byte for byte and every kept store row for row, the attestation ledger unchanged, each derived head
equal to the head before, the chat log rebuilt from raw alone, and no embedder built. Two resets
equal one; a fresh integration adds no causal edge. The rebuilt atom table holds exactly the atoms
the rebuilt rows name. On a lazy instance with a backlog the reset builds no embedder and the
backlog stays. A purged transcript is skipped and the reset finishes. The dry run writes nothing,
on an instance at rest not even a file's modification time, and names each store. A live run, a
held lock, a changed geometry, a pruned ledgered commit and a derived store outside the data
directory each refuse with every byte unchanged and no marker written. A failure injected at each
of the nine steps leaves the marker, the daemon and the repository sync refuse, and a second reset
converges; a held version purged before the rerun is skipped by it, a changed embedder while the
reset is pending is refused with the way out, and every script that writes refuses while the
marker stands. A destroy's refusals each write nothing, a root holding a repository, a foreign file
or another instance among them, and its dry run leaves every file's size and modification time as
they were. Run as commands in processes of their own, on a root at rest whose catalog has an empty
\code{-wal} beside it, the dry runs of destroy, reset and fill change no file, and on a root that
does not exist they create none. A stream part linked to the catalog's file, by a symbolic or a
hard link, or to a note, is refused by the destroy, the fill and the cache. With a dialogue turn, a
model's reply, an inbox item and a curated item added, a confirmed destroy leaves exactly the open
digests, the atoms open items held, and the stream's downloaded files in the cache, and no note,
dialogue, inbox item, transcript or vector only a private item held. A new instance takes a cached vector with no embedder
call and the same values; under another key it embeds. A corrupted raw entry is refused and
removed. A purge removes the object and its atoms' vectors under two keys, and a later fill from an
instance still holding it does not restore them.

The landing log is checked on the same fixtures. A reset keeps every row of the log and each
chain's head hash, and so do two resets, a reset interrupted and finished, and a sync that heals a
stale index, whose one new row is its own landing. The steps that open the catalog are recorded at
every open and are exactly the two named. On a catalog with rows and no log, a reset without the
secret refuses, its dry run too, with every row unchanged and no log table created; with the
secret its dry run still writes nothing, and the reset writes the genesis at its first open, one
\code{store} row for each catalog row and nothing more, and every chain verifies. Under a named instance loaded through its configuration
overlay, the instance's own item is the one read and the default instance's does not serve; a
reset of a catalog that has its log needs no secret; a destroy needs none, names the item among
those it leaves, and no file of the filled cache holds a value of the log. A destroy's dry run on
an instance at rest, every store closed, creates and touches no file. Mutating any of the
mechanisms these tests name (one at a time, restoring each) makes at least one test fail.

\paragraph{The corpus at its root.} With the corpus overlay declaring its root, a read-only check
on the tree with the landing log in its base (SQLite opened \code{immutable=1}, the Lance table
only read, the catalog never opened through its class) resolved the root to the declared
directory \code{instances/corpus} under the home root, counted \corpusitems{} active catalog items,
\corpusunits{} current lexical units and \corpusrows{} vector rows, and answered one lexical query
(\emph{superconductivity phase transition}, top hit the Wikipedia article Superconductivity) in
\lexsec{}~s at a one-minute load average of \lexload{} on a shared machine, scored under the legacy
tokenizer: the index has no recorded tokenizer, so by the store's own rule it is the legacy one's,
and the next corpus landing will re-derive the whole index first. The corpus's catalog has no
landing log yet, so opening it through its class would write the log's genesis under the root, or
refuse without the corpus's secret after creating the log's tables. A snapshot of the whole root
before and after the check, \corpusentries{} entries by type, size and modification time, showed
none added, removed or changed. An earlier check opened the files with \code{mode=ro}, which on a
WAL database creates its \code{-wal} and \code{-shm} files; it created two such files under the
corpus root and changed no row.

\section{Limitations and future work}

A reset rebuilds current versions only, and keeps the catalog's rows rather than folding them from
the landing log. A replay of the log would rebuild every version landed since its chain began, and
the fold needs a catalog write outside its logging methods; both are follow-ups, and the versions
superseded before a catalog's genesis are in no ledger but the version chains. A reset excludes a
writer that starts after it and one that holds the supervisor lock, but not an on-demand writer
that takes no lock and started before it; closing that needs a lock every writer holds for its
whole run, a decision left open. The vector key
names the model but not the model file, so a requantized file under one name would share a key; a
measurement would say whether to add its digest. Code membership vectors do not enter the cache,
because no reader takes them back yet. The cost of a full reset of the corpus instance, and of the
lexical re-derivation its next landing will run, is not measured here, nor is the time of a fill
of the cache from the corpus, whose plan counts 1.0~GB of raw objects, 1.4~GB of source files and
8.2~GB of vectors. A destroy of the corpus refuses as its root stands: its tooling writes files at
the root that no namespace of the instance names, and whether they are the instance's state is
left open. A pruned ledgered commit
refuses every reset of its instance until a way to accept that loss is designed. The plane does not
record the embedder's backend, so a change of backend alone is not caught before a fill; catching
it needs the plane to record the backend beside its geometry, a schema change left open. The class
table is checked on fixtures; a store a lane writes that the fixtures do not exercise would escape it.

\section{Conclusion}

Treating every index as a view over kept bytes gives an instance a reset that loses nothing and a
destroy that leaves a shared, verifiable cache. The table that classifies stores is the one place
both operations read, so they mean the same thing in every instance.

\bibliographystyle{plainnat}
\bibliography{refs}

\end{document}
